\documentclass[12pt,a4paper]{article}
\usepackage[utf8]{inputenc}
\usepackage[T1]{fontenc}
\usepackage[portuguese]{babel}
\usepackage[left=3cm,top=3cm,right=2cm,bottom=2cm]{geometry}
\usepackage{times}
\usepackage{setspace}
\usepackage{indentfirst}
\usepackage{booktabs}
\usepackage{longtable}
\usepackage{url}
\usepackage{microtype}
\usepackage{amsmath}
\usepackage[colorlinks=true,linkcolor=black,urlcolor=blue,citecolor=black]{hyperref}

% Configurações de Espaçamento e Parágrafo ABNT
\onehalfspacing
\setlength{\parindent}{1.5cm}
\setlength{\parskip}{6pt}

% Atalho para grau Celsius
\newcommand{\celsius}{$^\circ$C}

\begin{document}

% --------------------------------------------------
% CAPA (ANONIMIZADA)
% --------------------------------------------------
\begin{titlepage}
    \begin{center}
        \textbf{UNIVERSIDADE FEDERAL [ANONIMIZADA]}\\
        \textbf{HOSPITAL UNIVERSITÁRIO [ANONIMIZADO]}\\
        \textbf{EDITAL N\textsuperscript{o} 01/2026 -- HU BRASIL}\\
        \textbf{PROGRAMA DE INICIAÇÃO CIENTÍFICA DA HU BRASIL (PIC/HU BRASIL)}\\
        \textbf{PROCESSO SELETIVO DE BOLSISTAS DE INICIAÇÃO CIENTÍFICA 2026-2027}
        
        \vspace{4cm}
        
        \Large\textbf{INVESTIGAÇÃO DE BACTÉRIAS MULTIRRESISTENTES DE IMPORTÂNCIA EPIDEMIOLÓGICA EM SUPERFÍCIES INANIMADAS HOSPITALARES DE CONTATO INDIRETO COMO POTENCIAIS RESERVATÓRIOS AMBIENTAIS}
        
        \vspace{4cm}
        
        \begin{flushright}
            \begin{minipage}{9.5cm}
                \begin{spacing}{1.1}
                    \textit{Projeto de Pesquisa (Anonimizado)} apresentado como requisito para inscrição no Processo Seletivo de Bolsas de Iniciação Científica do PIC/HU Brasil (Ciclo 2026-2027).\\
                    
                    \textbf{Área de conhecimento conforme CNPq:} Ciências da Saúde / Medicina / Doenças Infecciosas e Parasitárias
                \end{spacing}
            \end{minipage}
        \end{flushright}
        
        \vfill
        
        \textbf{[Localidade -- UF]}\\
        \textbf{2026}
    \end{center}
\end{titlepage}

\newpage
\pagenumbering{arabic} % Começa a contar a partir do sumário/resumo

% --------------------------------------------------
% 1. SUMÁRIO
% --------------------------------------------------
\tableofcontents
\newpage

% --------------------------------------------------
% 2. RESUMO
% --------------------------------------------------
\section{Resumo}
\noindent As infecções relacionadas à assistência à saúde (IRAS) constituem um grave problema de saúde pública, impactando diretamente a morbimortalidade de pacientes hospitalizados. Embora as superfícies de contato direto (como as mãos dos profissionais de saúde) recebam constante atenção em protocolos de higienização, superfícies de contato indireto e menor vigilância no ambiente hospitalar --- como cortinas divisórias de leitos, maçanetas das portas de acesso, suportes de soro e painéis de controle de bombas de infusão --- podem atuar como reservatórios persistentes de patógenos multirresistentes. Este estudo tem como objetivo avaliar a contaminação microbiológica e o perfil de susceptibilidade antimicrobiana de bactérias de importância clínica presentes nessas superfícies inanimadas secundárias no Hospital Universitário [ANONIMIZADO]. Trata-se de um estudo transversal descritivo no qual serão coletadas 100 amostras das quatro superfícies selecionadas em enfermarias ativas, utilizando swabs estéreis umedecidos em solução salina estéril. As amostras serão submetidas a enriquecimento em caldo BHI, semeadura em meios de cultura seletivos cromogênicos específicos para o isolamento de KPC/CRE, VRE e MRSA, seguidos de identificação por coloração de Gram e espectrometria de massas MALDI-TOF MS. O perfil de susceptibilidade aos antimicrobianos dos isolados clínicos será determinado pelo método de disco-difusão em ágar (Kirby-Bauer), em conformidade com as diretrizes vigentes do BrCAST/EUCAST. Como o estudo envolve exclusivamente superfícies físicas, inanimadas e coletivas de ambiente de livre circulação, não há coleta de material de seres humanos ou identificação de pacientes, dispensando a submissão prévia ao Comitê de Ética em Pesquisa (CEP), o que agiliza sua exequibilidade. Espera-se com este estudo mapear a carga microbiana de superfícies negligenciadas e fornecer subsídios científicos para a revisão e o aprimoramento das práticas de limpeza e desinfecção hospitalar, visando reduzir as taxas de transmissão cruzada.

\vspace{12pt}
\noindent \textbf{Palavras-chave:} Contaminação de Equipamentos; Serviço Hospitalar de Limpeza; Antibióticos; Infecção Hospitalar; Meio Ambiente de Instituições de Saúde.
\newpage

% --------------------------------------------------
% 3. INTRODUÇÃO
% --------------------------------------------------
\section{Introdução}

\subsection{Apresentação do tema}
O ambiente hospitalar é dinâmico e envolve a circulação contínua de indivíduos com diferentes estados de saúde, o que favorece a colonização e a dispersão de microrganismos patogênicos. A sobrevivência bacteriana em superfícies inanimadas próximas ao paciente tem sido amplamente documentada na literatura como um fator crítico para a ocorrência de surtos e para a endemicidade de infecções associadas aos cuidados de saúde. A dinâmica de transmissão cruzada envolve a transferência de patógenos das superfícies ambientais para as mãos dos profissionais e, subsequentemente, para os pacientes, ou diretamente das superfícies para dispositivos médicos invasivos.

Historicamente, o foco dos programas de controle de infecção e higienização hospitalar tem sido as superfícies de alta frequência de toque direto e os dispositivos médicos invasivos. Contudo, há uma ampla variedade de superfícies inanimadas que fazem parte da rotina assistencial e de internação que recebem menos vigilância microbiológica e menor frequência de desinfecção sistemática. Entre essas superfícies estão as cortinas de separação de leitos (frequentemente tocadas por pacientes e profissionais antes ou após o contato físico direto), as maçanetas das portas de acesso às enfermarias e quartos de isolamento, os suportes verticais de soro e os painéis de controle digitais de bombas de infusão. Tais elementos, devido ao seu posicionamento espacial estratégico no entorno imediato do leito, representam pontes de transmissão indireta e podem albergar patógenos por dias ou semanas.

A relevância científica de investigar esses reservatórios reside no fato de que bactérias multirresistentes de extrema relevância epidemiológica --- como as enterobactérias produtoras de carbapenemase do tipo KPC, enterococos resistentes à vancomicina (VRE) e \textit{Staphylococcus aureus} resistente à meticilina (MRSA) --- apresentam notável estabilidade ambiental e capacidade de persistir em superfícies hospitalares, atuando como fontes de surtos de difícil controle terapêutico. A determinação da contaminação ambiental por esses patógenos específicos no Hospital Universitário [ANONIMIZADO] preenche uma lacuna de vigilância epidemiológica ativa, saindo do rastreamento genérico e afunilando em alvos críticos para fundamentar ações de controle de infecção hospitalar direcionadas e eficazes.

\subsection{Definição do problema}
Diante do papel epidemiológico das superfícies ambientais na persistência de patógenos, levanta-se a seguinte questão de pesquisa:
\textbf{Superfícies inanimadas hospitalares de contato indireto e baixo monitoramento microbiológico atuam como reservatórios ativos para bactérias multirresistentes de relevância epidemiológica crítica (KPC, VRE e MRSA)?}

\subsection{Hipótese}
\begin{itemize}
    \item \textbf{Hipótese Geral:} Superfícies inanimadas secundárias de contato frequente no entorno do paciente (cortinas divisórias, maçanetas das portas de acesso, suportes de soro e painéis de bombas de infusão), que não passam por desinfecção concorrente sistemática, apresentam contaminação significativa e atuam como reservatórios de bactérias multirresistentes sob vigilância epidemiológica (KPC, VRE e MRSA).
    \item \textbf{Hipóteses Secundárias:}
    \begin{enumerate}
        \item A presença de patógenos multirresistentes específicos difere significativamente entre as superfícies têxteis (cortinas), que propiciam maior aderência e sobrevivência de cocos Gram-positivos (como VRE e MRSA), e as superfícies poliméricas ou metálicas (maçanetas, suportes e painéis), propícias a bacilos Gram-negativos multirresistentes (como CRE/KPC).
        \item A taxa de contaminação e a diversidade de microrganismos resistentes são maiores em superfícies localizadas na vizinhança imediata do paciente (maçanetas de portas e cortinas divisórias), em comparação com superfícies de contato intermitente (suportes de soro e painéis de bombas).
        \item A contaminação ambiental por cepas de MRSA, VRE e enterobacterias produtoras de carbapenemase (KPC) correlaciona-se com a alta pressão de seleção de antimicrobianos nas enfermarias ativas da instituição.
    \end{enumerate}
\end{itemize}
\newpage

% --------------------------------------------------
% 4. OBJETIVO GERAL e 5. OBJETIVOS ESPECÍFICOS
% --------------------------------------------------
\section{Objetivo geral}
Avaliar o perfil de contaminação ambiental e a presença de bactérias multirresistentes de importância epidemiológica crítica (especificamente enterobactérias produtoras de carbapenemase do tipo KPC, enterococos resistentes à vancomicina - VRE, e \textit{Staphylococcus aureus} resistente à meticilina - MRSA) em superfícies inanimadas de contato indireto e sob baixo monitoramento em enfermarias do Hospital Universitário [ANONIMIZADO], estimando o seu papel potencial como reservatórios ambientais de transmissão cruzada.

\section{Objetivos específicos}
\begin{enumerate}
    \item Rastrear e isolar especificamente bactérias multirresistentes (MRSA, VRE e CRE/KPC) a partir de amostras de swabs coletados em quatro tipos de superfícies inanimadas: cortinas de separação de leitos, maçanetas das portas de acesso, suportes de soro e painéis de controle de bombas de infusão.
    \item Confirmar a identificação taxonômica das cepas multirresistentes isoladas por meio de coloração de Gram e espectrometria de massas MALDI-TOF MS.
    \item Caracterizar fenotipicamente o perfil de susceptibilidade antimicrobiana e confirmar os mecanismos de resistência dos isolados bacterianos (como resistência à meticilina por cefoxitina, resistência à vancomicina e detecção fenotípica de carbapenemases) seguindo as diretrizes do BrCAST.
    \item Avaliar comparativamente as taxas de contaminação por patógenos multirresistentes específicos entre as superfícies amostradas, identificando as áreas críticas que exigem maior rigor nos protocolos de higienização.
    \item Propor e formular recomendações práticas ao Serviço de Controle de Infecção Hospitalar (SCIH) baseadas na prevalência local das bactérias multirresistentes detectadas, visando o aprimoramento dos protocolos de desinfecção concorrente e terminal.
\end{enumerate}
\newpage

% --------------------------------------------------
% 6. JUSTIFICATIVA
% --------------------------------------------------
\section{Justificativa}
As infecções relacionadas à assistência à saúde (IRAS) impõem encargos significativos aos sistemas de saúde através do aumento no tempo de internação dos pacientes, morbimortalidade elevada e expressivo incremento nos custos operacionais com terapias antimicrobianas de largo espectro. O controle dessas infecções exige uma abordagem multifacetada que vai além da higiene das mãos dos profissionais, englobando a qualidade microbiológica do ambiente assistencial imediato. A investigação científica direcionada a superfícies de contato indireto e monitoramento escasso reveste-se de grande relevância prática, pois locais como cortinas, suportes de soro, maçanetas de portas e painéis de equipamentos eletrônicos frequentemente escapam dos processos regulares e rigorosos de desinfecção hospitalar mecânica.

As cortinas de privacidade hospitalares, por exemplo, são feitas majoritariamente de tecidos que absorvem umidade e matéria orgânica, fornecendo condições físicas favoráveis para a adesão e viabilidade bacteriana de longo prazo de patógenos como VRE e MRSA. Simultaneamente, maçanetas de portas, suportes de soro e painéis de equipamentos eletrônicos são manipulados de forma repetitiva por profissionais, acompanhantes e pelos próprios pacientes, funcionando como fômites ideais para a dispersão de bacilos Gram-negativos multirresistentes produtores de carbapenemases (KPC/CRE). O rastreamento específico desses microrganismos de alta relevância epidemiológica no Hospital Universitário [ANONIMIZADO] gerará dados cruciais para que o Serviço de Controle de Infecção Hospitalar redesenhe as estratégias de higienização, isolamento de contato e desinfecção terminal.

Sob a ótica teórica, este projeto fornece dados locais sobre a circulação e perfil de susceptibilidade de bactérias nosográficas e sua persistência sob condições de estresse ambiental. A opção por focar a pesquisa exclusivamente em superfícies físicas e amostras do ambiente inanimado (sem o envolvimento direto de seres humanos, sem coleta de dados clínicos individuais ou manipulação direta de pacientes) elimina os riscos bioéticos associados ao manuseio de dados pessoais sensíveis e procedimentos invasivos em indivíduos hospitalizados. Isso viabiliza a execução célere das etapas do projeto sem a necessidade de tramitação longa no Comitê de Ética em Pesquisa (CEP), otimizando o cumprimento rigoroso do cronograma de 12 meses da Iniciação Científica e garantindo a aplicabilidade imediata dos resultados laboratoriais nas rotinas assistenciais do hospital.
\newpage

% --------------------------------------------------
% 7. REVISÃO DA LITERATURA
% --------------------------------------------------
\section{Revisão da literatura}
A transmissão de microrganismos multirresistentes a partir de reservatórios ambientais constitui uma via reconhecida de aquisição de infecções hospitalares. Estudos pioneiros demonstraram que diversos patógenos nosocomiais conseguem sobreviver por longos períodos em superfícies inanimadas secas. \textit{Staphylococcus aureus} (incluindo MRSA) e \textit{Acinetobacter baumannii} podem persistir em superfícies de plástico ou metal por meses, enquanto enterococos sobrevivem por semanas em materiais têxteis.

O estudo clássico realizado por OHL et al. (2012) [PMID 22464039] investigou longitudinalmente a dinâmica de contaminação de cortinas de privacidade hospitalares. Os autores constataram que as cortinas novas ou higienizadas tornavam-se rapidamente contaminadas com bactérias patogênicas em um curto intervalo de tempo após a sua colocação nos quartos dos pacientes. Em poucos dias de uso contínuo, patógenos oportunistas como \textit{Staphylococcus aureus} e enterococos resistentes à vancomicina (VRE) foram isolados dos tecidos das cortinas, evidenciando a rapidez com que essas barreiras físicas tornam-se reservatórios microbiológicos e destacando a ineficiência de trocas ou higienizações esporádicas.

Corroborando a preocupação com a persistência de patógenos em ambientes críticos, SHEK et al. (2017) [PMID 28413115] conduziram um estudo transversal em uma unidade de queimados e cirurgia plástica para avaliar as cortinas de privacidade. Os autores observaram taxas extremamente elevadas de contaminação bacteriana nestas superfícies, o que é alarmante dada a alta vulnerabilidade de pacientes queimados a infecções bacterianas invasivas. O estudo confirmou que superfícies manipuladas de forma repetida que não entram nos protocolos de desinfecção diária sistemática acumulam carga patogênica expressiva, configurando riscos graves de transmissão cruzada para pacientes que já apresentam a barreira cutânea comprometida.

Para mitigar a contaminação desses reservatórios têxteis, SCHWEIZER et al. (2021) [PMID 34750366] conduziram um ensaio clínico randomizado piloto avaliando diferentes métodos de desinfecção química ativa para reduzir a carga bacteriana em cortinas de privacidade. A pesquisa avaliou a eficácia de sprays desinfetantes e lenços umedecidos com agentes biocidas aplicados diretamente nas cortinas em uso. Os resultados mostraram que intervenções direcionadas podem reduzir temporariamente a carga bacteriana, contudo, a recontaminação ocorre rapidamente se a higiene das mãos dos profissionais e a frequência de limpeza geral não forem rigorosamente mantidas. Esse achado ressalta a complexidade da manutenção da esterilidade relativa e reforça a necessidade de vigilância constante do ambiente.

Adicionalmente, FERREIRA et al. (2016) [PMID 27289247] avaliaram a contaminação por MRSA em quartos de pacientes colonizados ou infectados por essa bactéria no Brasil. A pesquisa analisou várias superfícies físicas, demonstrando correlação direta entre a presença do paciente portador e o isolamento de cepas geneticamente idênticas de MRSA nas superfícies do leito, maçanetas de portas e mesas de cabeceira e maçanetas. O estudo evidenciou que a contaminação ambiental não é homogênea e que superfícies de alto contato atuam como os principais epicentros microbiológicos no quarto do paciente, exigindo processos específicos de desinfecção terminal e concorrente para quebrar o ciclo de transmissão no ambiente hospitalar brasileiro.

No contexto brasileiro, a disseminação de enterobactérias produtoras de carbapenemase (\textit{Klebsiella pneumoniae} carbapenemase - KPC) e de enterococos resistentes à vancomicina (VRE) tem se consolidado como um desafio de controle epidemiológico. SILVA et al. (2021) [DOI: 10.33448/rsd-v10i11.19565] identificaram altas taxas de colonização de superfícies inanimadas hospitalares por enterobactérias com perfil de resistência a carbapenêmicos, ressaltando o papel crítico do ambiente como reservatório silencioso. Adicionalmente, surtos recentes de VRE em hospitais terciários brasileiros documentados por SILVA et al. (2024) [DOI: 10.1590/S2237-96222024v34e20240135.en] evidenciam que a contaminação cruzada por meio de superfícies de contato indireto e pias é um vetor central na persistência e alastramento de clones multirresistentes, tornando a vigilância ambiental ativa destas superfícies um elemento estratégico de controle sanitário.
\newpage

% --------------------------------------------------
% 8. ESTADO DA TÉCNICA
% --------------------------------------------------
\section{Estado da Técnica}
O monitoramento da higiene ambiental e da carga microbiológica em serviços de saúde evoluiu significativamente, dispondo atualmente de abordagens visuais, químicas e microbiológicas. Cada metodologia apresenta vantagens metodológicas e limitações operacionais distintas na prática de controle de infecção hospitalar.

A inspeção visual constitui o método mais simples e tradicional de avaliação da limpeza, porém é altamente subjetiva e incapaz de detectar contaminação invisível a olho nu ou de avaliar a presença de patógenos viáveis. Métodos químicos, como o uso de marcadores fluorescentes aplicados em superfícies antes da higienização para posterior leitura com luz ultravioleta, são úteis para avaliar a adesão dos profissionais de limpeza aos protocolos físicos, mas não medem a carga bacteriana residual. A bioluminescência de ATP (trifosfato de adenosina) mede a presença de matéria orgânica total (microbiana e biológica humana) em poucos segundos, fornecendo um indicador rápido de sujeira em superfícies. Contudo, o teste de ATP não discrimina entre detritos biológicos viáveis ou inviáveis e não identifica as espécies bacterianas presentes, tampouco seus respectivos perfis de resistência aos antibióticos.

Nesse panorama, a utilização de \textbf{meios de cultura cromogênicos seletivos} representa o estado da técnica mais avançado e específico para a vigilância microbiológica ativa de microrganismos multirresistentes. Ao contrário do cultivo convencional em meios tradicionais --- que recupera uma grande diversidade de bactérias da microbiota comum e de menor importância epidemiológica (tornando o processo laborioso e pouco direcionado) ---, os meios cromogênicos contêm inibidores e substratos enzimáticos específicos. Isso permite que bactérias como MRSA, VRE e enterobactérias produtoras de carbapenemase (KPC/CRE) cresçam seletivamente com colorações distintas e características em 24 a 48 horas. O cultivo seletivo, seguido pelo teste de susceptibilidade aos antimicrobianos (TSA) pelo método de disco-difusão (Kirby-Bauer), permanece como o \textbf{padrão-ouro} para caracterização fenotípica e confirmação de mecanismos de resistência regulados pelo BrCAST/EUCAST, permitindo ações rápidas de controle de infecção hospitalar.
\newpage

% --------------------------------------------------
% 9. METODOLOGIA
% --------------------------------------------------
\section{Metodologia}

\subsection{Delineamento do estudo e local da pesquisa}
Este é um estudo observacional transversal com abordagem descritiva e caráter analítico in vitro, focado no isolamento e na identificação microbiológica de bactérias a partir de superfícies inanimadas. O estudo será conduzido nas enfermarias clínicas e cirúrgicas adultas do Hospital Universitário [ANONIMIZADO]. As análises laboratoriais de cultivo, identificação bioquímica e testes de susceptibilidade antimicrobiana serão executadas nas dependências do Laboratório de Microbiologia [ANONIMIZADO] da própria instituição.

\subsection{Superfícies amostradas e critérios de seleção}
Serão coletadas 100 amostras microbiológicas de superfícies ambientais, distribuídas de forma homogênea ($n = 25$ para cada tipo de superfície):
\begin{enumerate}
    \item \textbf{Cortinas de separação de leitos:} Amostragem da zona de toque frequente (borda lateral anterior da cortina, a uma altura média de 1,2 a 1,5 metros do piso).
    \item \textbf{Maçanetas de portas de acesso:} Coleta na superfície de contato manual das maçanetas das portas de acesso às enfermarias e aos quartos de isolamento.
    \item \textbf{Suportes de soro:} Amostragem da haste vertical metálica ajustável na área de pegada manual dos profissionais de saúde.
    \item \textbf{Painéis de controle de bombas de infusão:} Coleta na superfície do teclado numérico e botões de acionamento do dispositivo de infusão ativa no leito.
\end{enumerate}
\noindent\textbf{Critérios de Inclusão:} Superfícies pertencentes a leitos de enfermaria ocupados por pacientes com tempo de internação superior a 48 horas.\\
\textbf{Critérios de Exclusão:} Superfícies localizadas em leitos desocupados ou que passaram por processo documentado de desinfecção terminal ou concorrente há menos de 2 horas da coleta programada.

\subsection{Coleta e transporte das amostras}
As coletas serão realizadas semanalmente por um período de três meses (entre o 3\textsuperscript{o} e o 5\textsuperscript{o} mês do projeto). Para cada superfície, será utilizado um swab de algodão estéril previamente umedecido em solução salina estéril a 0,85\%. Para superfícies planas delimitadas (cortinas, maçanetas de portas e suportes de soro), a coleta será padronizada friccionando-se o swab em movimentos firmes unidirecionais e rotatórios sobre uma área de $10\times10$ cm ($100\text{ cm}^2$), delimitada temporariamente por um gabarito de plástico estéril descartável. Para o painel da bomba de infusão, devido ao formato irregular, o swab será friccionado uniformemente sobre toda a extensão dos botões e área de toque do teclado.

Imediatamente após a fricção, os swabs serão acondicionados individualmente em tubos com meio de transporte Stuart identificados com códigos de anonimização (ex. CRT-01, MAC-01, SUP-01, BOM-01). Os tubos serão mantidos sob refrigeração em caixas térmicas (4\celsius{} a 8\celsius{}) e transportados ao laboratório para processamento em até duas horas após a coleta.

\subsection{Processamento laboratorial e isolamento bacteriano}
No laboratório, os swabs serão transferidos para tubos contendo 5~mL de Caldo Infusão Cérebro Coração (BHI --- \textit{Brain Heart Infusion}) com o objetivo de realizar o enriquecimento primário da amostra e aumentar a viabilidade microbiana e sensibilidade de recuperação. Os tubos serão incubados em estufa microbiológica a 35\celsius{} ($\pm$ 2\celsius{}) por 24 horas sob condições aeróbicas.

Após o período de incubação em caldo BHI, alíquotas de 10~$\mu$L da cultura enriquecida serão semeadas por esgotamento (técnica de estrias isoladas) em placas de Petri contendo meios de cultura seletivos cromogênicos específicos para o isolamento dos alvos multirresistentes delineados:
\begin{itemize}
    \item \textbf{Ágar Seletivo Cromogênico para KPC/CRE (ex: CHROMagar KPC ou similar):} Destinado ao isolamento seletivo de bacilos Gram-negativos resistentes a carbapenêmicos. Colônias com colorações características (azul-metálico para espécies do grupo KPC ou rosa para outras enterobactérias resistentes) serão consideradas suspeitas.
    \item \textbf{Ágar Seletivo Cromogênico para MRSA (ex: CHROMagar MRSA ou similar):} Utilizado para detecção seletiva de \textit{Staphylococcus aureus} resistente à meticilina. O crescimento de colônias de cor rosa-intensa/mauve indicará forte suspeita de MRSA.
    \item \textbf{Ágar Seletivo Cromogênico para VRE (ex: CHROMagar VRE ou similar):} Empregado para o isolamento seletivo de enterococos resistentes à vancomicina. Colônias de coloração azul ou azul-esverdeada serão indicativas de VRE (\textit{Enterococcus faecium}/\textit{Enterococcus faecalis}).
    \item \textbf{Ágar MacConkey:} Mantido como meio de controle e isolamento secundário para bacilos Gram-negativos de contaminação geral.
\end{itemize}
As placas serão incubadas a 35\celsius{} por 24 a 48 horas sob condições aeróbicas, realizando-se a leitura periódica das colônias de coloração característica para cada microrganismo sob vigilância.

\subsection{Identificação taxonômica por MALDI-TOF MS e Coloração de Gram}
A identificação das colônias isoladas nos meios seletivos cromogênicos será realizada por espectrometria de massas, otimizando o tempo e a precisão do diagnóstico taxonômico:
\begin{enumerate}
    \item \textbf{Triagem Inicial (Coloração de Gram):} Confirmação morfotintorial por coloração de Gram a partir de colônias suspeitas isoladas nos meios cromogênicos, permitindo a verificação preliminar da morfologia bacteriana (cocos Gram-positivos ou bacilos Gram-negativos) e pureza das culturas.
    \item \textbf{Identificação por Espectrometria de Massas MALDI-TOF MS:} As colônias isoladas serão submetidas à identificação em nível de gênero e espécie pela tecnologia MALDI-TOF MS (\textit{Matrix-Assisted Laser Desorption/Ionization Time-of-Flight Mass Spectrometry}). Pequenas frações das colônias puras serão transferidas para placas metálicas de aço (placa de alvo), recobertas com uma matriz orgânica apropriada (geralmente ácido alfa-ciano-4-hidroxicinâmico - HCCA) e analisadas no espectrômetro de massas. Os perfis espectrais de proteínas obtidos (principalmente proteínas ribossomais) serão comparados eletronicamente com o banco de dados do fabricante do equipamento (ex. MALDI Biotyper), gerando scores de identificação confiáveis para a espécie bacteriana analisada.
\end{enumerate}

\subsection{Teste de Sensibilidade aos Antimicrobianos (TSA)}
O perfil de susceptibilidade e a confirmação das resistências de importância epidemiológica serão determinados pelo método de disco-difusão em ágar Mueller-Hinton (Kirby-Bauer), seguindo rigorosamente os diâmetros de halos e pontos de corte regulamentados pelo BrCAST/EUCAST vigente.
\begin{itemize}
    \item \textbf{Confirmação de MRSA:} Isolados de \textit{Staphylococcus aureus} que apresentarem diâmetro de halo do disco de Cefoxitina (30~$\mu$g) inferior ao ponto de corte de resistência serão confirmados como MRSA. Também serão testados discos de Penicilina G (1~UI/10~$\mu$g), Eritromicina (15~$\mu$g), Clindamicina (2~$\mu$g), Sulfametoxazol-Trimetoprima (1,25/23,75~$\mu$g), Ciprofloxacina (5~$\mu$g), Linezolida (30~$\mu$g) e Teicoplanina (30~$\mu$g).
    \item \textbf{Confirmação de VRE:} Confirmada pela resistência aos discos de Vancomicina (5~$\mu$g) e Teicoplanina (30~$\mu$g). Serão testados também Ampicilina (2~$\mu$g), Linezolida (30~$\mu$g), Gentamicina de alta carga (120~$\mu$g) e Estreptomicina de alta carga (300~$\mu$g).
    \item \textbf{Confirmação de CRE/KPC e Não-fermentadores Resistentes a Carbapenêmicos:} Confirmada pela resistência aos carbapenêmicos Imipenem (10~$\mu$g), Meropenem (10~$\mu$g) e Ertapenem (10~$\mu$g). Adicionalmente, serão testados discos de Amoxicilina/Ácido Clavulânico (20/10~$\mu$g), Ceftazidima (30~$\mu$g), Cefepima (30~$\mu$g), Ceftazidima-Avibactam (30/20~$\mu$g), Ciprofloxacina (5~$\mu$g), Amicacina (30~$\mu$g), Sulfametoxazol-Trimetoprima e Polimixina B (utilizando método de microdiluição em caldo se exigido pelo comitê para confirmação de perfil de pan-resistência).
    \item \textbf{Controle de Qualidade:} Serão utilizadas cepas de referência internacionais para controle de qualidade dos discos e do ágar Mueller-Hinton em todos os lotes testados (\textit{Staphylococcus aureus} ATCC 25923, \textit{Escherichia coli} ATCC 25922 e \textit{Pseudomonas aeruginosa} ATCC 27853).
\end{itemize}

\subsection{Análise de dados e tratamento estatístico}
Os dados microbiológicos e as variáveis das superfícies (tipo de superfície, local da coleta) serão organizados em banco de dados eletrônico no programa Microsoft Excel. A análise estatística descritiva consistirá no cálculo de frequências absolutas e relativas das superfícies contaminadas e das espécies bacterianas identificadas, bem como dos percentuais de sensibilidade e resistência aos antimicrobianos.

A análise inferencial será conduzida para testar a associação entre o tipo de superfície e a presença de patógenos específicos ou bactérias multirresistentes. Para tanto, serão aplicados os testes estatísticos Qui-quadrado de Pearson ou o Teste Exato de Fisher, adotando-se o nível de significância estatística de $p < 0,05$. As análises estatísticas serão efetuadas por meio do software de distribuição livre R (versão atualizada).

\subsection{Aspectos éticos e legais}
A presente proposta de pesquisa caracteriza-se como um estudo in vitro de vigilância ambiental no qual as amostras microbiológicas são isoladas exclusivamente a partir de superfícies inanimadas coletivas e públicas instaladas no ambiente físico hospitalar (maçanetas de portas, hastes de soro, painéis eletrônicos e cortinas). Não haverá qualquer coleta de fluidos biológicos humanos, aplicação de questionários, análise de prontuários clínicos, ou procedimentos de intervenção direta ou indireta em pacientes, acompanhantes ou profissionais de saúde.

Assim, de acordo com as diretrizes estabelecidas pelas Resoluções n\textsuperscript{o} 466/2012 e n\textsuperscript{o} 510/2016 do Conselho Nacional de Saúde (CNS), que dispõem sobre pesquisas envolvendo seres humanos no Brasil, o presente projeto não apresenta elementos constitutivos que configurem pesquisa com seres humanos. Por conseguinte, este estudo não exige submissão ou tramitação para análise ética no Comitê de Ética em Pesquisa (CEP)/CONEP, permitindo a execução imediata das coletas de amostras a partir do início da vigência das atividades do projeto, desde que autorizada formalmente pela Diretoria de Ensino e Pesquisa do Hospital Universitário [ANONIMIZADO] por meio da obtenção da Carta de Anuência Institucional (conforme Modelo em Anexo).
\newpage

% --------------------------------------------------
% 10. RESULTADOS ESPERADOS
% --------------------------------------------------
\section{Resultados esperados}
Com a execução deste estudo, espera-se alcançar os seguintes resultados técnico-científicos:
\begin{enumerate}
    \item \textbf{Mapeamento de Patógenos Multirresistentes:} Identificar as taxas de contaminação específicas por MRSA, VRE e CRE/KPC nas superfícies hospitalares de contato indireto, evidenciando quais fômites atuam como principais reservatórios desses microrganismos de vigilância.
    \item \textbf{Relação Epidemiológica de Riscos:} Compreender a distribuição espacial e a taxa de circulação ambiental dessas bactérias multirresistentes nas enfermarias ativas do Hospital Universitário [ANONIMIZADO].
    \item \textbf{Caracterização de Perfis de Resistência Críticos:} Determinar os padrões de susceptibilidade cruzada a outros antimicrobianos nas cepas de MRSA, VRE e KPC isoladas, indicando o grau de multirresistência e a presença de possíveis isolados com perfil de sensibilidade reduzida a novos compostos de uso restrito (como Ceftazidima-Avibactam).
    \item \textbf{Subsídio a Protocolos de Barreira e Higiene:} Elaborar um relatório técnico analítico detalhado e direcionado à Comissão de Controle de Infecção Hospitalar (CCIH/SCIH) e ao Serviço de Higiene e Limpeza. Estes dados empíricos servirão de base para fundamentar a implantação de precauções de contato adicionais, revisão dos cronogramas de troca/lavagem de cortinas de leito e especificação dos saneantes ideais para desinfecção dos painéis eletrônicos e maçanetas.
    \item \textbf{Formação de Recursos Humanos:} Desenvolver as competências científicas e laboratoriais do estudante bolsista na área de microbiologia de multirresistentes e epidemiologia hospitalar in vitro, capacitando-o para o trabalho em segurança do paciente.
\end{enumerate}
\newpage

% --------------------------------------------------
% 11. CRONOGRAMA DO PROJETO
% --------------------------------------------------
\section{Cronograma do Projeto}
A execução do projeto de pesquisa terá duração de 12 meses (coincidente com a vigência da bolsa do PIC/HU Brasil, de 01/09/2026 a 31/08/2027), estruturada conforme o detalhamento mensal abaixo:

\begin{spacing}{1.0}
\begin{longtable}{lcccccccccccc}
\caption{Cronograma de execução das atividades do projeto.}\\
\toprule
\textbf{Etapa / Atividade} & \textbf{M1} & \textbf{M2} & \textbf{M3} & \textbf{M4} & \textbf{M5} & \textbf{M6} & \textbf{M7} & \textbf{M8} & \textbf{M9} & \textbf{M10} & \textbf{M11} & \textbf{M12} \\
\midrule
\endfirsthead
\multicolumn{13}{c}{Sinalização de continuação da tabela.}\\
\toprule
\textbf{Etapa / Atividade} & \textbf{M1} & \textbf{M2} & \textbf{M3} & \textbf{M4} & \textbf{M5} & \textbf{M6} & \textbf{M7} & \textbf{M8} & \textbf{M9} & \textbf{M10} & \textbf{M11} & \textbf{M12} \\
\midrule
\endhead
\bottomrule
\endfoot
1. Revisão e busca bibliográfica & X & X & & & & & & & & & & \\
2. Obtenção de autorizações locais & X & & & & & & & & & & & \\
3. Planejamento e treino de lab & & X & & & & & & & & & & \\
4. Coleta semanal nas enfermarias & & & X & X & X & & & & & & & \\
5. Isolamento em meios cromogênicos & & & X & X & X & X & & & & & & \\
6. Confirmação e identificação bioquímica & & & & & & X & X & X & & & & \\
7. Teste de Sensibilidade e confirmação (TSA) & & & & & & & X & X & X & & & \\
8. Tabulação e estatística & & & & & & & & & X & X & & \\
9. Relatório Parcial do PIC & & & & & & & & & & X & & \\
10. Discussão dos dados com CCIH & & & & & & & & & & & X & \\
11. Relatório Final de IC e artigo & & & & & & & & & & & X & X \\
12. Submissão ao Congresso de IC & & & & & & & & & & & & X \\
\bottomrule
\end{longtable}
\end{spacing}
\newpage

% --------------------------------------------------
% 12. ORÇAMENTO DO PROJETO
% --------------------------------------------------
\section{Orçamento do Projeto}
Os recursos necessários para a realização deste projeto de pesquisa são compostos por materiais de consumo específicos para a atividade laboratorial de baixo custo. A infraestrutura permanente (estufas, autoclave, microscópios, capela de fluxo laminar) está disponível nas dependências do Hospital Universitário, não demandando gastos adicionais.

\begin{spacing}{1.0}
\begin{longtable}{clcrr}
\caption{Orçamento detalhado dos materiais de consumo do projeto.}\\
\toprule
\textbf{Item} & \textbf{Descrição do Insumo / Especificação} & \textbf{Unid.} & \textbf{Qtd.} & \textbf{Valor Total (R\$)} \\
\midrule
\endfirsthead
\toprule
\textbf{Item} & \textbf{Descrição do Insumo / Especificação} & \textbf{Unid.} & \textbf{Qtd.} & \textbf{Valor Total (R\$)} \\
\midrule
\endhead
\bottomrule
\endfoot
1 & Swabs estéreis de algodão para coleta c/ transporte & Caixa & 1 & 150,00 \\
2 & Caldo Infusão Cérebro Coração (BHI) frasco 500g & Frasco & 1 & 420,00 \\
3 & Ágar Cromogênico Seletivo p/ KPC/CRE desidratado & Frasco & 1 & 680,00 \\
4 & Ágar Cromogênico Seletivo p/ MRSA desidratado & Frasco & 1 & 640,00 \\
5 & Ágar Cromogênico Seletivo p/ VRE desidratado & Frasco & 1 & 650,00 \\
6 & Ágar MacConkey desidratado - frasco 500g & Frasco & 1 & 240,00 \\
7 & Ágar Mueller-Hinton desidratado - frasco 500g & Frasco & 1 & 350,00 \\
8 & Discos de antibióticos p/ TSA (Kirby-Bauer) & Frasco & 8 & 360,00 \\
9 & Kit de reagentes para coloração de Gram & Kit & 1 & 160,00 \\
10 & Placas de Petri plásticas descartáveis 90x15mm & Caixa & 1 & 290,00 \\
11 & Alças bacteriológicas descartáveis estéreis 10~$\mu$L & Pacote & 3 & 120,00 \\
12 & EPIs (luvas de látex, máscaras, toucas) & Pacote & 4 & 240,00 \\
13 & Tubos de ensaio de vidro autoclaváveis 15 mL & Pacote & 1 & 180,00 \\
14 & Fita indicadora de esterilização autoclavável (rolo) & Rolo & 2 & 60,00 \\
\midrule
\multicolumn{4}{r}{\textbf{VALOR TOTAL GERAL DOS MATERIAIS}} & \textbf{R\$ 4.500,00} \\
\bottomrule
\end{longtable}
\end{spacing}
\newpage

% --------------------------------------------------
% 13. PLANO DE TRABALHO DO ALUNO
% --------------------------------------------------
\section{Plano de trabalho do(a) aluno(a)}

\subsection{Descrição das atividades do(a) aluno(a) (12 meses)}
O(A) estudante bolsista desempenhará papel ativo em todas as fases do projeto de pesquisa, atuando sob supervisão direta do orientador e dos técnicos do laboratório. Suas atividades incluem:
\begin{itemize}
    \item \textbf{Mês 1 (Set/2026):} Estudo de manuais de biossegurança laboratorial; revisão da literatura focando em contaminação ambiental e resistência bacteriana; redação do fichamento dos artigos principais.
    \item \textbf{Mês 2 (Out/2026):} Auxílio na submissão de documentos e obtenção da anuência formal do hospital; treinamento técnico prático no laboratório de microbiologia, cobrindo o preparo de meios de cultura, esterilização (autoclavagem) e técnicas de semeadura microbiológica.
    \item \textbf{Mês 3 a 5 (Nov/2026 a Jan/2027):} Realização sistemática e semanal das coletas de swabs nas enfermarias utilizando técnica asséptica e gabarito; transporte térmico das amostras e inoculação inicial em caldo BHI no laboratório.
    \item \textbf{Mês 6 a 8 (Fev/2027 a Abr/2027):} Isolamento bacteriano em meios cromogênicos seletivos (MRSA, VRE, KPC); realização de coloração de Gram; identificação automatizada por espectrometria de massas (MALDI-TOF MS) e coloração de Gram.
    \item \textbf{Mês 7 a 9 (Mar/2027 a Mai/2027):} Preparo de suspensão bacteriana na escala 0,5 de McFarland e semeadura em placas Mueller-Hinton; aplicação de discos de antimicrobianos para confirmação do perfil de multirresistência por TSA; leitura e medição milimétrica dos halos de inibição.
    \item \textbf{Mês 9 e 10 (Mai/2027 a Jun/2027):} Tabulação dos resultados em planilha eletrônica; aplicação de testes estatísticos de associação (Qui-quadrado/Fisher) com auxílio do orientador; redação do relatório parcial das atividades de Iniciação Científica para envio à coordenação do PIC.
    \item \textbf{Mês 11 (Jul/2027):} Discussão dos perfis de contaminação com a equipe de controle de infecção (CCIH) e elaboração de material visual (como panfletos ou cartazes explicativos) sobre a higienização de superfícies acessórias.
    \item \textbf{Mês 12 (Ago/2027):} Redação do relatório final de pesquisa da Iniciação Científica, contendo os resultados e a discussão confrontada com a literatura científica atualizada; preparação de resumo científico para submissão a congresso acadêmico da instituição ou de microbiologia nacional.
\end{itemize}

\subsection{Cronograma das atividades do(a) aluno(a)}
O cronograma de atividades do bolsista coincide com o do projeto global, cobrindo os 12 meses de vigência da bolsa do PIC/HU Brasil (01/09/2026 a 31/08/2027):

\begin{spacing}{1.0}
\begin{longtable}{lcccccccccccc}
\caption{Cronograma de atividades do estudante bolsista.}\\
\toprule
\textbf{Etapa / Atividade do Estudante} & \textbf{M1} & \textbf{M2} & \textbf{M3} & \textbf{M4} & \textbf{M5} & \textbf{M6} & \textbf{M7} & \textbf{M8} & \textbf{M9} & \textbf{M10} & \textbf{M11} & \textbf{M12} \\
\midrule
\endfirsthead
\toprule
\textbf{Etapa / Atividade do Estudante} & \textbf{M1} & \textbf{M2} & \textbf{M3} & \textbf{M4} & \textbf{M5} & \textbf{M6} & \textbf{M7} & \textbf{M8} & \textbf{M9} & \textbf{M10} & \textbf{M11} & \textbf{M12} \\
\midrule
\endhead
\bottomrule
\endfoot
1. Revisão e biossegurança & X & X & & & & & & & & & & \\
2. Treinamento prático e anuência & & X & & & & & & & & & & \\
3. Coleta de swab nas enfermarias & & & X & X & X & & & & & & & \\
4. Inoculação BHI e semeadura cromogênica & & & X & X & X & X & & & & & & \\
5. Identificação por MALDI-TOF MS e Gram & & & & & & X & X & X & & & & \\
6. Realização de Testes (TSA) & & & & & & & X & X & X & & & \\
7. Organização e tratamento de dados & & & & & & & & & X & X & & \\
8. Relatório Parcial do PIC & & & & & & & & & & X & & \\
9. Interações com CCIH e material & & & & & & & & & & & X & \\
10. Relatório Final e artigo científico & & & & & & & & & & & X & X \\
\bottomrule
\end{longtable}
\end{spacing}
\newpage

% --------------------------------------------------
% 14. BIBLIOGRAFIA
% --------------------------------------------------
\section{Bibliografia}
\begin{spacing}{1.1}
\noindent FERREIRA, Adriano Menis et al. Contamination of environmental surfaces by methicillin-resistant Staphylococcus aureus (MRSA) in rooms of inpatients with MRSA-positive body sites. \textbf{Brazilian Journal of Microbiology}, v. 47, n. 3, p. 610-616, 2016. DOI: 10.1016/j.bjm.2016.04.009. (PMID: 27289247).

\vspace{12pt}
\noindent OHL, Michael et al. Hospital privacy curtains are frequently and rapidly contaminated with potentially pathogenic bacteria. \textbf{American Journal of Infection Control}, v. 40, n. 10, p. 904-906, 2012. DOI: 10.1016/j.ajic.2011.12.017. (PMID: 22464039).

\vspace{12pt}
\noindent SCHWEIZER, Marin L. et al. Assessment of cleaning methods on bacterial burden of hospital privacy curtains: a pilot randomized controlled trial. \textbf{Scientific Reports}, v. 11, n. 1, p. 21503, 2021. DOI: 10.1038/s41598-021-00936-4. (PMID: 34750366).

\vspace{12pt}
\noindent SHEK, Kevin et al. Rate of contamination of hospital privacy curtains on a burns and plastic surgery ward: a cross-sectional study. \textbf{Journal of Hospital Infection}, v. 96, n. 1, p. 54-58, 2017. DOI: 10.1016/j.jhin.2017.03.012. (PMID: 28413115).

\vspace{12pt}
\noindent ALMEIDA, Maria Tibiriçá Gonçalves de et al. Papel das superfícies inanimadas na transmissão de infecções relacionadas à assistência à saúde. \textbf{Revista de Epidemiologia e Controle de Infecção}, v. 8, n. 2, p. 140-145, 2018. DOI: 10.17058/reci.v8i2.10221.

\vspace{12pt}
\noindent BRASIL. Agência Nacional de Vigilância Sanitária (ANVISA). \textbf{Segurança do Paciente em Serviços de Saúde: Higienização das Mãos}. Brasília: Anvisa, 2023. Disponível em: \url{https://www.gov.br/anvisa/pt-br/assuntos/servicosdesaude/seguranca-do-paciente}. Acesso em: 25 jun. 2026.

\vspace{12pt}
\noindent BrCAST. \textbf{Tabela de pontos de corte para interpretação de testes de sensibilidade aos antimicrobianos}. Versão atualizada, 2026. Disponível em: \url{http://brcast.org.br/}. Acesso em: 25 jun. 2026.

\vspace{12pt}
\noindent WEINBERG, Jane et al. Survival of enterococci and staphylococci on hospital surfaces. \textbf{Journal of Clinical Microbiology}, v. 44, n. 2, p. 345-349, 2006. DOI: 10.1128/JCM.44.2.345-349.2006.

\vspace{12pt}
\noindent SILVA, Adriana Carla Silva de et al. Identificação fenotípica de bactérias da família Enterobacteriaceae com perfil de resistência em superfícies inanimadas e utensílios em Hospital Universitário. \textbf{Research, Society and Development}, v. 10, n. 11, e299101119565, 2021. DOI: 10.33448/rsd-v10i11.19565.

\vspace{12pt}
\noindent SILVA, Daniel Antunes Custódio da et al. Vancomycin-resistant Enterococcus infections in a hospital in Salvador, Bahia: a descriptive study, 2021-2023. \textbf{Epidemiologia e Serviços de Saúde}, v. 34, e20240135, 2024. DOI: 10.1590/S2237-96222024v34e20240135.en. (PMID: 40243735).
\end{spacing}
\newpage

% --------------------------------------------------
% 15. ANEXOS
% --------------------------------------------------
\section{Anexos}

\subsection{Modelo de Carta de Anuência Institucional para Pesquisa}

\begin{center}
    \textbf{CARTA DE ANUÊNCIA INSTITUCIONAL}
\end{center}

\noindent Eu, \_\_\_\_\_\_\_\_\_\_\_\_\_\_\_\_\_\_\_\_\_\_\_\_\_\_\_\_\_\_\_\_\_\_\_\_\_\_\_\_\_\_\_\_\_\_\_\_\_\_\_, na qualidade de Diretor(a) / Coordenador(a) do Serviço de Controle de Infecção Hospitalar / Direção de Ensino e Pesquisa do Hospital Universitário [ANONIMIZADO], declaro estar ciente dos termos da pesquisa intitulada \textbf{"Investigação de Bactérias Multirresistentes de Importância Epidemiológica em Superfícies Inanimadas Hospitalares de Contato Indireto como Potenciais Reservatórios Ambientais"}, sob responsabilidade da equipe científica anonimizada cadastrada sob o Edital n\textsuperscript{o} 01/2026 da HU Brasil.

Declaro que esta instituição autoriza a coleta de dados microbiológicos por meio de swabs friccionados em superfícies inanimadas públicas de uso coletivo (maçanetas de portas de acesso, cortinas de leitos, suportes de soro e painéis de controle de bombas de infusão), em enfermarias sob jurisdição técnica desta administração.

A autorização é concedida sob a ciência de que:
\begin{enumerate}
    \item Não haverá coleta de materiais biológicos de pacientes, acompanhantes ou profissionais, tampouco acesso a dados identificadores individuais dos mesmos.
    \item As atividades serão conduzidas sem interferência com o atendimento assistencial rotineiro ou com as normas de limpeza concorrente e terminal do hospital.
    \item Os resultados obtidos serão compartilhados com os setores técnicos competentes do hospital (CCIH e Gerência de Higiene) para fins de vigilância sanitária local e aprimoramento interno das práticas de segurança do paciente.
\end{enumerate}

\vspace{1cm}
\noindent Hospital Universitário [ANONIMIZADO], \_\_\_\_ de \_\_\_\_\_\_\_\_\_\_\_\_\_\_\_\_ de 2026.

\vspace{1.5cm}
\begin{center}
    ________________________________________________________\\
    Assinatura e carimbo do Responsável Institucional
\end{center}

\newpage
\subsection{Ficha de Coleta de Amostras de Superfícies}

\begin{center}
    \textbf{FICHA DE REGISTRO E RASTREABILIDADE MICROBIANA}
\end{center}

\noindent \textbf{Identificador da Porta (Código):} \_\_\_\_\_\_\_\_\_\_\_\_\_\_\_\_\_\_\_\_ \hfill \textbf{Enfermaria / Setor:} \_\_\_\_\_\_\_\_\_\_\_\_\_\_\_\_\_\_\_\_ \\
\textbf{Data da Coleta:} \_\_\_\_\_/\_\_\_\_\_/2026 \hfill \textbf{Hora da Coleta:} \_\_\_\_\_:\_\_\_\_\_ \\
\textbf{Tempo decorrido desde a última higienização concorrente (horas):} \_\_\_\_\_\_\_\_

\vspace{12pt}
\noindent \textbf{Checklist de Coleta de Amostras (Marcar após acondicionamento no meio Stuart):}

\begin{spacing}{1.0}
\begin{longtable}{|l|l|c|p{6.5cm}|}
\hline
\textbf{Cód. Amostra} & \textbf{Superfície} & \textbf{Efetuada (S/N)} & \textbf{Observações físicas da superfície} \\
\hline
CRT-\_\_\_\_\_\_ & Cortina Divisória & & \\
\hline
MAC-\_\_\_\_\_\_ & Maçaneta da Porta & & \\
\hline
SUP-\_\_\_\_\_\_ & Suporte de Soro & & \\
\hline
BOM-\_\_\_\_\_\_ & Painel da Bomba & & \\
\hline
\end{longtable}
\end{spacing}

\vspace{12pt}
\noindent \textbf{Coletor Responsável (Rubrica):} \_\_\_\_\_\_\_\_\_\_\_\_\_\_\_\_\_\_\_\_ \\
\textbf{Data de Recebimento no Laboratório:} \_\_\_\_\_/\_\_\_\_\_/2026 \hfill \textbf{Hora:} \_\_\_\_\_:\_\_\_\_\_ \\
\textbf{Técnico de Laboratório Responsável (Assinatura):} \_\_\_\_\_\_\_\_\_\_\_\_\_\_\_\_\_\_\_\_

\end{document}
